\documentclass{minmal}
\begin{document}
\begin{center}
\textbf{prøve om varmevekslere:}
\end{center}
\textbf{Oppg. 1:}\nl[1]
\noindent Forklar hva som mens med begrepene:\nl[1]
a) Spesifikk varmekapasitet\nl[1]
b) Spesifikk smeltevarme\nl[1]
c) Spesifikk fordampingsvarme\nl[1]
d) Hvilken verdi har vannets spesifikke varmekapasitet?\nl[1]
e) Hvilken verdi har vannets spesifikke fordampingsvarme? \nl[2]
\\
\textbf{Oppg. 2:}

\noindent a) 8,3kg vann som har en temperatur på 17\textsuperscript{0}C skal varmes til 73\textsuperscript{0}C. Beregn energimengden som trengs for å få til dette når du vet at vannets spesifikke varmekapasitet er C\textsubscript{P,vann}=4,2kJ/kg\textsuperscript{0}C. \nl[1]
b) Hvor mye energi kreves for å fordampe 6,3kg vann ved 100\textsuperscript{0}C når du vet at vannets spesifikke fordampingsvarme er l\textsubscript{f} =2256kJ/kg? \nl[1]
c) Hvor mye energi trengs for å smelte 6,3kg is når du vet at vannets spesifikke smeltevarme er l\textsubscript{s}=334kJ/kg.\nl[1] 
d) Hvor mye energi trengs for å fordampe 7,8kg is som har en temperatur på -13\textsuperscript{0}C?\nl[1]
\newpage
\noindent\textbf{Oppg. 3:}\nl[1]
Gitt følgende varmeveksler:\nl[1]
\varmeveksler{75}{55}{?}{7}{23}{25}\nl[1]
a) Hvor mye energi overføres i denne varmeveksleren?\nl[1]
b) Hva blir temperaturen ut av varmeveksleren på kald side ($T_{k,ut}$)?\newpage
Gitt følgende varmeveksler:
\varmeveksler{66}{46}{37}{20}{?}{19}
c)	Bestem den overførte varmemengden i varmeveksleren.\nl[1]
d)	Bestem $F_v$. \newpage
\textbf{Formler:}\nl[1]
\begin{figure}[h]
    \centering
    \includegraphics[width=0.8\linewidth]{formel_VV.png}
    
    \label{fig:placeholder}
\end{figure}\\
\textbf{ Fordampingsvarme:}

 \noindent $Q_f=m\cdot l_f$\\
$l_f=2256kJ/kg$\\

\noindent\textbf{Smeltevarme:}\\

 \noindent
 $Q_s=m\cdot l_s$\\
\noindent
$l_s=334kJ/kg$\\
  m=masse (kg)\\
\textit{C\textsubscript{p,vann}=4,2kJ/kgC\\  
C\textsubscript{p,is}=2,1kJ/kgC}


\end{document}